  
\addcontentsline{toc}{chapter}{\twkkai{对话录·其二}}  
\markboth{\twkkai{对话录·其二}}{}  
\pagecolor{black}
\color{white}
\quad\quad \twkkai{\lettrine[,lraise=0.2,lhang=0.5]{\textbf{猫}}{：}是时候醒醒了，我的朋友。

狄利克雷：……好啦。又被你吵醒了。有什么收获吗？

猫：没什么大不了，只是又在街上随便逛了逛。不过我再次听到了人群嘈杂的噪音中传来的动听的人造音乐。似乎和你的研究有点联系呢。

狄利克雷：你有录下来吗？

猫：当然。我这枚奖章就是CD；我的爪子可以当唱针。

狄利克雷：你已经说过了。别大惊小怪——它原本只是怕你走丢而已……

猫：可别这么说，我才不是你的宠物。这是我自己得到的。

狄利克雷：你可别忘恩负义；这奖章是我留给你找到使者的——可不是给你闹着玩的。

猫：好，好……不过你真的不想听听录下来的音乐吗？是韦伯作曲的Op.73第二乐章。

狄利克雷：嗯……你播来听听。

猫：好。

\noindent\adjustbox{center}{\includegraphics[scale=0.7]{Motive-A.pdf}}

猫：气氛真是紧张\marginnote{
\centering
{\Large ♪}

\vspace{0.4em}

\qraudio{Motive-A}{\MotiveASound}{https://raw.githubusercontent.com/kelvinlamresearchmath-web/finding-dirichlet-site/master/Motive-A.mp3}

\vspace{0.3em}
}[-4.1cm]
{\small ♪}——开头\kern-15pt\lily[hpadding=0pt]{
    \override Staff.StaffSymbol.color       = #white
    \override Staff.Clef.color              = #white
    \override Staff.TimeSignature.color     = #white
    \override Staff.KeySignature.color      = #white
    \override Staff.BarLine.color           = #white
    \override Staff.LedgerLineSpanner.color = #white
    \override NoteHead.color                = #white
    \override Stem.color                    = #white
    \override Beam.color                    = #white
    \override Flag.color                    = #white
    \override Rest.color                    = #white
    \override Accidental.color              = #white
    \override Dots.color                    = #white
    \override Slur.color                    = #white
    \override Tie.color                     = #white
    \override Script.color                  = #white
    \override DynamicText.color             = #white
    \override Hairpin.color                 = #white
    \override TextScript.color              = #white
    \override TupletBracket.color           = #white
    \override TupletNumber.color            = #white
  \clef treble
  \key d \major
  \omit Staff.TimeSignature
  f16-. a( d') a(
  f') d'( a') f'( 
  d'') a'( f'') d''-.
  a''-\tenuto( e'' f'' d'') \bar "||"
}\kern5pt 实在是错落有致；这便是一个由同一和声音型反复换序进行排列变换构成的动机展开，最后加上音符“E”来收束变形——构成一个d小三和弦（上行）琶音。

\rule{0pt}{5pt}

\noindent\adjustbox{center}{\includegraphics[scale=0.7]{Motive-B.pdf}}

猫：啊\marginnote{
\centering
{\Large ♪}

\vspace{0.4em}

\qraudio{Motive-B}{\MotiveBSound}{https://raw.githubusercontent.com/kelvinlamresearchmath-web/finding-dirichlet-site/master/Motive-B.mp3}

\vspace{0.3em}
}[-3.8cm]
{\small ♪}……最紧张的部分——\kern-10pt\lily[hpadding=0pt]{
    \override Staff.StaffSymbol.color       = #white
    \override Staff.Clef.color              = #white
    \override Staff.TimeSignature.color     = #white
    \override Staff.KeySignature.color      = #white
    \override Staff.BarLine.color           = #white
    \override Staff.LedgerLineSpanner.color = #white
    \override NoteHead.color                = #white
    \override Stem.color                    = #white
    \override Beam.color                    = #white
    \override Flag.color                    = #white
    \override Rest.color                    = #white
    \override Accidental.color              = #white
    \override Dots.color                    = #white
    \override Slur.color                    = #white
    \override Tie.color                     = #white
    \override Script.color                  = #white
    \override DynamicText.color             = #white
    \override Hairpin.color                 = #white
    \override TextScript.color              = #white
    \override TupletBracket.color           = #white
    \override TupletNumber.color            = #white
  \clef treble
  \key d \major
  \omit Staff.TimeSignature
  r16 c'''-\tenuto( bes'' a'' g'' fis'' a'' fis'')
r16 g'-\tenuto( bes' d'' fis'' g'' a'' bes'')
\bar "||"
}\kern5pt 和紧接着张力更强的第二句\kern-10pt\lily[hpadding=0pt]{
    \override Staff.StaffSymbol.color       = #white
    \override Staff.Clef.color              = #white
    \override Staff.TimeSignature.color     = #white
    \override Staff.KeySignature.color      = #white
    \override Staff.BarLine.color           = #white
    \override Staff.LedgerLineSpanner.color = #white
    \override NoteHead.color                = #white
    \override Stem.color                    = #white
    \override Beam.color                    = #white
    \override Flag.color                    = #white
    \override Rest.color                    = #white
    \override Accidental.color              = #white
    \override Dots.color                    = #white
    \override Slur.color                    = #white
    \override Tie.color                     = #white
    \override Script.color                  = #white
    \override DynamicText.color             = #white
    \override Hairpin.color                 = #white
    \override TextScript.color              = #white
    \override TupletBracket.color           = #white
    \override TupletNumber.color            = #white
  \clef treble
  \key d \major
  \omit Staff.TimeSignature
  r16 ees''' (-\tenuto d''' c'''! bes'' a'' g'' fis'')
r16 bes'-\tenuto( d'' g'' bes'' c'''! cis''' d''')
\bar "||"
}\kern5pt ——已然到来。我的朋友，你听——它们不仅表层结构一致——都是休止、起音、同一音型；而且深层结构都是前半部分有一个下行型的启动，后半部分有一个分解和弦式的展开；而第二句动机是第一句动机的规整版强化。两者共同构成了一个复合动机。

\rule{0pt}{5pt}

\noindent\adjustbox{center}{\includegraphics[scale=0.7]{Motive-C.pdf}}

猫：再听\marginnote{
\centering
{\Large ♪}

\vspace{0.4em}

\qraudio{Motive-C}{\MotiveCSound}{https://raw.githubusercontent.com/kelvinlamresearchmath-web/finding-dirichlet-site/master/Motive-C.mp3}

\vspace{0.3em}
}[-4.3cm]
{\small ♪}！最后一小节\kern-15pt\lily[hpadding=0pt]{
    \override Staff.StaffSymbol.color       = #white
    \override Staff.Clef.color              = #white
    \override Staff.TimeSignature.color     = #white
    \override Staff.KeySignature.color      = #white
    \override Staff.BarLine.color           = #white
    \override Staff.LedgerLineSpanner.color = #white
    \override NoteHead.color                = #white
    \override Stem.color                    = #white
    \override Beam.color                    = #white
    \override Flag.color                    = #white
    \override Rest.color                    = #white
    \override Accidental.color              = #white
    \override Dots.color                    = #white
    \override Slur.color                    = #white
    \override Tie.color                     = #white
    \override Script.color                  = #white
    \override DynamicText.color             = #white
    \override Hairpin.color                 = #white
    \override TextScript.color              = #white
    \override TupletBracket.color           = #white
    \override TupletNumber.color            = #white
    \clef treble
    \key d \major
    \omit Staff.TimeSignature
    f!16-\tenuto( a) d'-. f'!-.
    a'-. d''-. f''!-. a''-.
    d'''4 r4^\fermata
    \bar "||"
}\kern5pt 尾部尖锐的高音“D”，配合那无尽的延长休止符该有多美啊——婉转空灵。不过，我的朋友——听我分析这么久，你还想知道这一切和你研究的有什么关系吗？

狄利克雷：不妨说说。

猫：你不觉得，它们讨论的都是变中的不变性\index{bian@变中的不变性}吗？由几个音符组成的动机，经过那么多的变换——移调、逆行、倒影、扩大、紧缩\index{dong@动机（motif）!移调、逆行、倒影、扩大、紧缩}等等——竟然还能保持其主题风格不变。尤其是刚才带延长符号的最后一小节最能体现——它全部由“F”“A”“D”三个元素构成，但四个一组排列之后就变成了有层次感的又一个d小三和弦（上行）琶音了。

狄利克雷：变中的不变性……的确值得斟酌。和我1837年发现的定理很像呢。

猫：正是正是。

守门人：\begin{gather*}
    f\left(\frac{az+b}{cz+d}\right)=(cz+d)^k f(z),\quad g\left(\frac{az+b}{cz+d}\right)=(cz+d)^k g(z).
    \end{gather*}

猫：……又来了。他守门太累了——又在打瞌睡说梦话了。不过，他说的话似乎和我刚才所说的有点联系。你看，后来的数学家——我们暂且称之为现代数学家——发现了一种新的变换形式。它对应的正是守门人口中所说的一种变换：\(z\mapsto \frac{az+b}{cz+d}\)。这种变换叫莫比乌斯变换\index{mobiwusi@莫比乌斯变换}；满足这种变换的函数就叫作模形式\index{moxing@模形式}。

狄利克雷：看来我作为一个旧时代的数学家是不是错过了一点什么……

猫：不！不不……千万不能这么讲，我的朋友。我只是感慨一下。满足这个映射，就能像动机那般，每一个点\(z\)的坐标都可以被旋转、放缩、位移；但是它——我是指满足这个映射的模形式——它始终保持全纯性\index{moxing@模形式!全纯性}、群结构(\(ad-bc=1\))\index{moxing@模形式!群结构（$ad-bc=1$）}、权（或阶）\(k\)不变\index{moxing@模形式!权 $k$ 不变}。

狄利克雷：嗯。我能体会到你话里那种令人惊奇与赞叹的本质事物。

猫：你还应该体会到刚才听到的乐句和你研究的领域的联系。

狄利克雷：你是说——刚才我们所感受到的音乐的情绪一致性和不变性，可以对应……

猫：喵！是的——那种不变的、贯穿始终的紧张情绪，可以和莫比乌斯变换对应；无论乐句中的动机怎么变化，传达给听众的依旧是揪心的感觉。

狄利克雷：真是奥妙无穷啊。不过，我派你完成的任务恐怕你是完成不了了。

猫：为什么这么说？

狄利克雷：从你最近老是丢失奖章一事来看，我有一种不祥的预感……

猫：你这是对我的不信任呐，我的朋友。

狄利克雷：并无此意。只是……你先留在这吧。

\rule{0pt}{20pt}

\noindent\adjustbox{center}{\includegraphics[scale=1]{trebleclef.preview.pdf}}

}






